\section{Reduction and the bounded cohomology detector}
\label{sec:reduction-detector}

We first separate the conformal reduction from the quantitative part of
the proof. We then construct a differential form on a simplex of
probability vectors that detects simplicial volume. This form will remain fixed while
the metric and the probability map are deformed.

\subsection{The conformal reduction}

We use the following result, stated here in the precise form needed.

\begin{theorem}[OpenAI, Corollary 1.5 of
  \cite{OpenAIScalarInessentiality2026}]
\label{thm:vanishing-input}
Let \(M\) be a closed connected oriented smooth manifold of positive
dimension. If \(M\) admits a smooth Riemannian metric with nonnegative
scalar curvature, then its real simplicial volume is zero.
\end{theorem}

In particular, the hypothesis permits scalar curvature to vanish
identically. Theorem~\ref{thm:vanishing-input} is the only result from
the companion manuscript invoked below without proof.

The Yamabe reformulation of the scalar curvature inequality is
standard; \cite[Lemma 2.1]{MinZhengZhu2026} records the equivalence of
its uniform pointwise and integral forms. We include the short
classwise reduction that identifies the use of the vanishing input.

\begin{proposition}
\label{prop:yamabe-reduction}
Fix \(n\ge3\). Suppose that a constant \(C_n<\infty\) has the following
property: every closed connected oriented \(n\)-manifold with a smooth
metric \(h\) of constant scalar curvature \(-1\) satisfies
\begin{equation}
  \|M\|\le C_n\Vol_h(M).
  \label{eq:scalar-minus-one-target}
\end{equation}
Then every smooth Riemannian metric \(g\) on every such manifold
satisfies
\begin{equation}
  \|M\|\le C_n\int_M(\Scal_g^-)^{n/2}\,dV_g .
  \label{eq:yamabe-transfer}
\end{equation}
\end{proposition}

\begin{proof}
The assertion is immediate when \(\|M\|=0\), so assume that
\(\|M\|>0\). Write \(Y(M,[g])\) for the Yamabe constant of the
conformal class of \(g\):
\[
Y(M,[g])=
\inf_{u\in C^\infty(M),\,u>0}
\frac{\displaystyle\int_M
  \left(\frac{4(n-1)}{n-2}|du|_g^2+\Scal_g u^2\right)dV_g}
{\displaystyle\left(\int_Mu^{2n/(n-2)}\,dV_g\right)^{(n-2)/n}}.
\]
The Yamabe theorem, developed through the work of Yamabe, Trudinger,
Aubin, and Schoen \cite{Yamabe1960,Trudinger1968,Aubin1976,Schoen1984},
gives a conformal metric attaining this value and having constant
scalar curvature; see also \cite[Section 11]{LeeParker1987}.
Its scalar curvature is negative by
Theorem~\ref{thm:vanishing-input}. After a constant rescaling, it becomes a
metric \(h\) with \(\Scal_h=-1\). Evaluating the Yamabe quotient at the
constant function for this minimizing metric gives
\[
  -Y(M,[g])=\Vol_h(M)^{2/n}.
\]
For each positive \(u\), H\"older's inequality gives
\[
\int_M\Scal_g u^2\,dV_g
\ge
-\left(\int_M(\Scal_g^-)^{n/2}\,dV_g\right)^{2/n}
 \left(\int_Mu^{2n/(n-2)}\,dV_g\right)^{(n-2)/n}.
\]
The gradient term in the Yamabe quotient is nonnegative. Taking the
infimum therefore yields
\[
-Y(M,[g])
\le \left(\int_M(\Scal_g^-)^{n/2}\,dV_g\right)^{2/n}.
\]
Consequently \(\Vol_h(M)\le\int_M(\Scal_g^-)^{n/2}\,dV_g\), and
\eqref{eq:scalar-minus-one-target} implies
\eqref{eq:yamabe-transfer}.
\end{proof}

For the rest of the proof we fix \(M^n\) and \(g\) with
\(\Scal_g=-1\), and prove \eqref{eq:scalar-minus-one-target}. We may
continue to assume \(\|M\|>0\).

\subsection{A bounded cocycle and its probability form}

Let \(\Gamma=\pi_1(M)\). This group is countable because \(M\) has a
finite triangulation. We use two standard facts about bounded
cohomology with trivial real coefficients. First, duality between the
\(\ell^1\)-seminorm on homology and the supremum seminorm on bounded
cohomology gives
\begin{equation}
\|M\|=
\sup\left\{\langle\beta,[M]_{\R}\rangle:
 \beta\in H_b^n(M;\R),\ \|\beta\|_\infty\le1\right\}.
\label{eq:bounded-duality}
\end{equation}
See \cite[Section 1.1]{Gromov1982}.
Second, the isometric Mapping Theorem identifies \(H_b^k(M;\R)\) with
\(H_b^k(\Gamma;\R)\) for every \(k\ge2\), through the classifying map.
The version used here applies to connected countable CW spaces and
preserves the seminorms. We use the Mapping Theorem of Gromov
\cite[Section 3.1]{Gromov1982} and its canonical group formulation in
\cite[Theorems 8.3 and 8.6]{IvanovNotes}.

We use the homogeneous model for group cochains. Thus an \(n\)-cochain
is a bounded function on \(\Gamma^{n+1}\), invariant under simultaneous
left translation, and its coboundary is the alternating sum obtained by
omitting one argument. Alternation is a cochain projection of norm at
most one and is equivariantly cochain homotopic to the identity through
maps bounded in each degree. It therefore preserves bounded cohomology
classes \cite[Section~4.10, Proposition~4.26]{Frigerio2016}.
By \eqref{eq:bounded-duality}, the Mapping Theorem, and
approximation of the seminorm by cocycle representatives, there is a
homogeneous alternating cocycle
\[
  c:\Gamma^{n+1}\longrightarrow\R,\qquad \|c\|_\infty\le1,
\]
whose corresponding bounded singular cohomology class has value
\begin{equation}
  I:=\langle[c],[M]_{\R}\rangle\ge\frac12\|M\|.
  \label{eq:detector-value}
\end{equation}
For completeness, one can choose a class in
\eqref{eq:bounded-duality} with value at least \(3\|M\|/4\), choose a
group representative of norm at most \(1+\varepsilon\), alternate it,
and divide by \(1+\varepsilon\), with \(0<\varepsilon\le1/2\). Replacing
the cocycle by its negative first, if necessary, fixes the sign.

We allow a finite number of copies of every group label. More
precisely, let \(\mathcal L=\Gamma\times\{1,\ldots,s\}\), with the free
left \(\Gamma\)-action on the first factor, and evaluate \(c\) on
\(\mathcal L^{n+1}\) by projection to \(\Gamma^{n+1}\). The same
notation will be used when further finite copies of \(\mathcal L\) are
introduced. Extend \(c\) multilinearly to finitely supported vectors
on \(\mathcal L\).

Let
\[
\Delta_{\mathcal L}=
\left\{p=(p_\ell)_{\ell\in\mathcal L}:
 p_\ell\ge0,\ \sum_\ell p_\ell=1,\ p\text{ has finite support}\right\}.
\]
On each finite-dimensional face define the \(n\)-form
\begin{equation}
\eta_p(v_1,\ldots,v_n)=n!\,c(p,v_1,\ldots,v_n).
\label{eq:probability-form}
\end{equation}
These definitions agree on common faces.

We call a map \(F=(F_\ell)\) with nonnegative coordinates and
\(\sum_\ell F_\ell^2=1\) a \emph{root probability map}. For the maps
used here, whose coordinate supports are locally finite, the squared
coordinates \(p_\ell=F_\ell^2\) give the associated map to
\(\Delta_{\mathcal L}\).

\begin{proposition}
\label{prop:detector}
The form \(\eta\) is closed on every finite-dimensional face of
\(\Delta_{\mathcal L}\). There is a smooth \(\Gamma\)-equivariant map
\[
 F:\widetilde M\longrightarrow\ell^2(\mathcal L),\qquad
 F_\ell\ge0,\qquad \sum_\ell F_\ell^2=1,
\]
with finitely many label orbits and locally finite compact supports for
its coordinates. If \(p_\ell=F_\ell^2\), then the descended form
satisfies
\begin{equation}
  \int_M p^*\eta=I.
  \label{eq:detector-integral}
\end{equation}
\end{proposition}

\begin{proof}
Write \(\epsilon(v)=\sum_\ell v_\ell\) for the mass of a finitely
supported vector. Multilinear extension of the cocycle equation gives
\[
0=\sum_{j=0}^{n+1}(-1)^j\epsilon(w_j)
 c(w_0,\ldots,\widehat w_j,\ldots,w_{n+1}).
\]
Choose \(w_0\) of mass one and let \(w_1,\ldots,w_{n+1}\) have mass
zero. It follows that \(c(w_1,\ldots,w_{n+1})=0\). All tangent vectors
to a probability simplex have mass zero. Differentiating
\eqref{eq:probability-form} now shows \(d\eta=0\).

To construct \(F\), choose finitely many evenly covered open sets in
\(M\) and smooth nonnegative bumps subordinate to them whose nonzero
sets cover \(M\). Lift each bump to one chosen sheet and take all its
deck translates, denoting the resulting functions by \(b_\ell\).
Their supports are compact and locally finite on \(\widetilde M\).
Set
\[
 F_\ell=\frac{b_\ell}{\left(\sum_j b_j^2\right)^{1/2}}.
\]
The denominator is positive, and the stated properties follow.

Choose a finite oriented triangulation of \(M\). On its lift, assign
one label to every vertex equivariantly, and let \(p^{\mathrm{aff}}\)
be the probability map affine in barycentric coordinates on each
simplex. The homotopy
\((1-t)p+t p^{\mathrm{aff}}\) is equivariant and has locally finite
support. Applying Stokes' theorem to each simplex times \([0,1]\), and
cancelling the common faces, shows that its pullback integral is
constant. On a simplex with ordered vertex labels
\(\ell_0,\ldots,\ell_n\), write
\[
p^{\mathrm{aff}}=e_{\ell_0}
+\sum_{j=1}^nt_j(e_{\ell_j}-e_{\ell_0}).
\]
Alternation makes the value of
\(c(p^{\mathrm{aff}},dp^{\mathrm{aff}},\ldots,dp^{\mathrm{aff}})\)
on the coordinate frame constant and equal to
\(c(\ell_0,\ldots,\ell_n)\). The coordinate simplex has volume
\(1/n!\), so the factor \(n!\) in \eqref{eq:probability-form} makes its
integral exactly this cocycle value. Summing over the oriented
fundamental triangulation gives \eqref{eq:detector-integral}. The
identification with the chosen group class is the usual cochain
realization of the Mapping Theorem: an equivariant vertex labeling on
the universal cover evaluates the homogeneous group cochain, and two
such labelings give cochain-homotopic maps by the simplex prism
construction.
\end{proof}

\subsection{Measured equivariant families}

The deformations below use auxiliary probability choices. We record
their integration convention so that the detector value remains
unambiguous.

Let \((\Omega,\mu)\) be a standard Borel probability space with a
measure-preserving \(\Gamma\)-action. A measured equivariant family
over \(M\) consists of \(\Gamma\)-equivariant Borel data on
\(\widetilde M\times\Omega\), smooth in the \(\widetilde M\) variable.
Geometrically one may regard the diagonal quotient
\[
  \mathcal M=(\widetilde M\times\Omega)/\Gamma
\]
as a space whose leaves are ordinary smooth manifolds. If
\(\Gamma_\omega\) is the stabilizer of \(\omega\), the corresponding
leaf is \(\widetilde M/\Gamma_\omega\); the deck action is free and
proper in the \(\widetilde M\) factor.

Related equivariant measured bundles are central to the macroscopic
volume argument of Braun and Sauer \cite[Sections 1.2 and 4]{BraunSauer2021}.
Here we need only Borel transverse data and use the following direct
integration convention.

Choose a finite evenly covered
atlas \(U_a\) of \(M\), chosen lifts \(\widetilde U_a\), and a smooth
partition of unity \(\chi_a\) subordinate to the atlas. For an
integrable equivariant leafwise \(n\)-form \(\alpha\), set
\begin{equation}
\int_{\mathcal M}\alpha
 =
\sum_a\int_\Omega\int_{\widetilde U_a}
  (\chi_a\circ\pi)\,\alpha_\omega\,d\mu(\omega).
\label{eq:measured-integral}
\end{equation}
Equivariance and invariance of \(\mu\) show that this value is
independent of the lifts and partition. The same convention applies
to nonnegative densities. We write \(dV_h\) for the resulting measured
Riemannian density when \(h=h_\omega\) is a leafwise metric.

The data used in the first deformation will have bounds in a fixed
finite lifted atlas that are uniform over \(\omega\). These bounds
may depend on the instance: the original manifold and metric and the
finite sequence of constructions already made. For the forms used
here, these bounds ensure integrability over the compact base and
justify differentiation and Fubini. Leafwise Stokes' theorem follows from
\eqref{eq:measured-integral}: apply ordinary Stokes' theorem in each
lifted chart to the partitioned form; the terms containing
\(d\chi_a\) cancel on overlaps after the measure-preserving changes of
transverse coordinate.

In particular, suppose that two equivariant probability maps
\(p^0,p^1\) have uniformly locally finite coordinate supports and
uniform leafwise derivative bounds. Their affine homotopy has the
same properties. Stokes' theorem on the measured product with
\([0,1]\) gives
\[
\int_{\mathcal M}(p^0)^*\eta
 =
\int_{\mathcal M}(p^1)^*\eta.
\]
Thus every probability map constructed below continues to detect the
same number \(I\). Enlarging \(\Omega\) by an independent probability
space also preserves the integral by Fubini. We may discard invariant
null sets whenever needed. More general measured spaces of leaves
will enter the degree estimate in Section~\ref{sec:degree}; their chart and
integration conventions will be specified there.
